% Môn: Vật lý
% Lớp: 12
% Chương: Chương I. Vật lí nhiệt
% Bài: L12C1 Bài 6. Nhiệt hóa hơi riêng · Thực nghiệm xác định nhiệt nóng chảy riêng
% ===== Câu 87 =====
% ID: L12C1B6-196-TL
% Mức: NB
\dangbt{Thực nghiệm xác định nhiệt nóng chảy riêng}
\begin{bt}
% ID: L12C1B6-196-TL
% Mức: NB
Trình bày cách nhận biết và phương pháp giải dạng “Thực nghiệm xác định nhiệt nóng chảy riêng”. Phân tích nhận định sau và sửa lại nếu sai: “Mọi nhiệt lượng của dây nung đều chắc chắn chỉ dùng để nóng chảy mẫu.”
\loigiai{
Trước hết xác định hiện tượng, trạng thái và các đại lượng đã cho. Nêu định nghĩa hoặc hệ thức phù hợp, đổi về đơn vị SI, áp dụng đúng điều kiện và kiểm tra ý nghĩa kết quả. Nhận định cần đối chiếu: Có thể xác định λ từ năng lượng cung cấp và khối lượng nóng chảy do nguồn. Vì vậy nhận định đã cho sai; phát biểu đúng là: Có thể xác định λ từ năng lượng cung cấp và khối lượng nóng chảy do nguồn.
}
\end{bt}

% ===== Câu 91 =====
% ID: L12C1B6-197-TL
% Mức: NB
\begin{bt}
% ID: L12C1B6-197-TL
% Mức: NB
Trình bày cách nhận biết và phương pháp giải dạng “Thực nghiệm xác định nhiệt nóng chảy riêng”. Phân tích nhận định sau và sửa lại nếu sai: “Mọi nhiệt lượng của dây nung đều chắc chắn chỉ dùng để nóng chảy mẫu.”
\loigiai{
Trước hết xác định hiện tượng, trạng thái và các đại lượng đã cho. Nêu định nghĩa hoặc hệ thức phù hợp, đổi về đơn vị SI, áp dụng đúng điều kiện và kiểm tra ý nghĩa kết quả. Nhận định cần đối chiếu: Có thể xác định λ từ năng lượng cung cấp và khối lượng nóng chảy do nguồn. Vì vậy nhận định đã cho sai; phát biểu đúng là: Có thể xác định λ từ năng lượng cung cấp và khối lượng nóng chảy do nguồn.
}
\end{bt}

% ===== Câu 95 =====
% ID: L12C1B6-198-TL
% Mức: TH
\begin{bt}
% ID: L12C1B6-198-TL
% Mức: TH
Trình bày cách nhận biết và phương pháp giải dạng “Thực nghiệm xác định nhiệt nóng chảy riêng”. Phân tích nhận định sau và sửa lại nếu sai: “Có thể xác định λ từ năng lượng cung cấp và khối lượng nóng chảy do nguồn.”
\loigiai{
Trước hết xác định hiện tượng, trạng thái và các đại lượng đã cho. Nêu định nghĩa hoặc hệ thức phù hợp, đổi về đơn vị SI, áp dụng đúng điều kiện và kiểm tra ý nghĩa kết quả. Nhận định cần đối chiếu: Có thể xác định λ từ năng lượng cung cấp và khối lượng nóng chảy do nguồn. Vì vậy nhận định đã cho đúng.
}
\end{bt}

% ===== Câu 98 =====
% ID: L12C1B6-199-DS
% Mức: NB
\begin{ex}
% ID: L12C1B6-199-DS
% Mức: NB
Xét các phát biểu sau về dạng “Thực nghiệm xác định nhiệt nóng chảy riêng” (bộ số 4).
\choiceTF
{\True Có thể xác định λ từ năng lượng cung cấp và khối lượng nóng chảy do nguồn.}
{Mọi nhiệt lượng của dây nung đều chắc chắn chỉ dùng để nóng chảy mẫu.}
{\True Khi giải dạng “Thực nghiệm xác định nhiệt nóng chảy riêng”, cần xác định đúng trạng thái đầu, trạng thái cuối và quy ước dấu hoặc điều kiện áp dụng.}
{Có thể bỏ qua mọi dữ kiện về khối lượng, nhiệt độ và hiệu suất trong mọi bài toán thuộc dạng này.}
\loigiai{
\begin{itemchoice}
\item (Đúng) Có thể xác định λ từ năng lượng cung cấp và khối lượng nóng chảy do nguồn. (Vì): Có thể xác định λ từ năng lượng cung cấp và khối lượng nóng chảy do nguồn. b) S: Mọi nhiệt lượng của dây nung đều chắc chắn chỉ dùng để nóng chảy mẫu. c) Đ: Khi giải dạng “Thực nghiệm xác định nhiệt nóng chảy riêng”, cần xác định đúng trạng thái đầu, trạng thái cuối và quy ước dấu hoặc điều kiện áp dụng. d) S: Có thể bỏ qua mọi dữ kiện về khối lượng, nhiệt độ và hiệu suất trong mọi bài toán thuộc dạng này. Kết luận dựa trên định nghĩa, điều kiện áp dụng và quy ước trong SGK KNTT
\item (Sai) Mọi nhiệt lượng của dây nung đều chắc chắn chỉ dùng để nóng chảy mẫu.
\item (Đúng) Khi giải dạng “Thực nghiệm xác định nhiệt nóng chảy riêng”, cần xác định đúng trạng thái đầu, trạng thái cuối và quy ước dấu hoặc điều kiện áp dụng.
\item (Sai) Có thể bỏ qua mọi dữ kiện về khối lượng, nhiệt độ và hiệu suất trong mọi bài toán thuộc dạng này.
\end{itemchoice}
}
\end{ex}

% ===== Câu 99 =====
% ID: L12C1B6-200-TL
% Mức: VD
\begin{bt}
% ID: L12C1B6-200-TL
% Mức: VD
Trình bày cách nhận biết và phương pháp giải dạng “Thực nghiệm xác định nhiệt nóng chảy riêng”. Phân tích nhận định sau và sửa lại nếu sai: “Mọi nhiệt lượng của dây nung đều chắc chắn chỉ dùng để nóng chảy mẫu.”
\loigiai{
Trước hết xác định hiện tượng, trạng thái và các đại lượng đã cho. Nêu định nghĩa hoặc hệ thức phù hợp, đổi về đơn vị SI, áp dụng đúng điều kiện và kiểm tra ý nghĩa kết quả. Nhận định cần đối chiếu: Có thể xác định λ từ năng lượng cung cấp và khối lượng nóng chảy do nguồn. Vì vậy nhận định đã cho sai; phát biểu đúng là: Có thể xác định λ từ năng lượng cung cấp và khối lượng nóng chảy do nguồn.
}
\end{bt}

% ===== Câu 102 =====
% ID: L12C1B6-201-DS
% Mức: TH
\begin{ex}
% ID: L12C1B6-201-DS
% Mức: TH
Xét các phát biểu sau về dạng “Thực nghiệm xác định nhiệt nóng chảy riêng” (bộ số 1).
\choiceTF
{Mọi nhiệt lượng của dây nung đều chắc chắn chỉ dùng để nóng chảy mẫu.}
{\True Khi giải dạng “Thực nghiệm xác định nhiệt nóng chảy riêng”, cần xác định đúng trạng thái đầu, trạng thái cuối và quy ước dấu hoặc điều kiện áp dụng.}
{Có thể bỏ qua mọi dữ kiện về khối lượng, nhiệt độ và hiệu suất trong mọi bài toán thuộc dạng này.}
{\True Có thể xác định λ từ năng lượng cung cấp và khối lượng nóng chảy do nguồn.}
\loigiai{
\begin{itemchoice}
\item (Sai) Mọi nhiệt lượng của dây nung đều chắc chắn chỉ dùng để nóng chảy mẫu. (Vì): Mọi nhiệt lượng của dây nung đều chắc chắn chỉ dùng để nóng chảy mẫu. b) Đ: Khi giải dạng “Thực nghiệm xác định nhiệt nóng chảy riêng”, cần xác định đúng trạng thái đầu, trạng thái cuối và quy ước dấu hoặc điều kiện áp dụng. c) S: Có thể bỏ qua mọi dữ kiện về khối lượng, nhiệt độ và hiệu suất trong mọi bài toán thuộc dạng này. d) Đ: Có thể xác định λ từ năng lượng cung cấp và khối lượng nóng chảy do nguồn. Kết luận dựa trên định nghĩa, điều kiện áp dụng và quy ước trong SGK KNTT
\item (Đúng) Khi giải dạng “Thực nghiệm xác định nhiệt nóng chảy riêng”, cần xác định đúng trạng thái đầu, trạng thái cuối và quy ước dấu hoặc điều kiện áp dụng.
\item (Sai) Có thể bỏ qua mọi dữ kiện về khối lượng, nhiệt độ và hiệu suất trong mọi bài toán thuộc dạng này.
\item (Đúng) Có thể xác định λ từ năng lượng cung cấp và khối lượng nóng chảy do nguồn.
\end{itemchoice}
}
\end{ex}

% ===== Câu 103 =====
% ID: L12C1B6-202-TL
% Mức: VDC
\begin{bt}
% ID: L12C1B6-202-TL
% Mức: VDC
Trình bày cách nhận biết và phương pháp giải dạng “Thực nghiệm xác định nhiệt nóng chảy riêng”. Phân tích nhận định sau và sửa lại nếu sai: “Có thể xác định λ từ năng lượng cung cấp và khối lượng nóng chảy do nguồn.”
\loigiai{
Trước hết xác định hiện tượng, trạng thái và các đại lượng đã cho. Nêu định nghĩa hoặc hệ thức phù hợp, đổi về đơn vị SI, áp dụng đúng điều kiện và kiểm tra ý nghĩa kết quả. Nhận định cần đối chiếu: Có thể xác định λ từ năng lượng cung cấp và khối lượng nóng chảy do nguồn. Vì vậy nhận định đã cho đúng.
}
\end{bt}

% ===== Câu 106 =====
% ID: L12C1B6-203-DS
% Mức: TH
\begin{ex}
% ID: L12C1B6-203-DS
% Mức: TH
Xét các phát biểu sau về dạng “Thực nghiệm xác định nhiệt nóng chảy riêng” (bộ số 5).
\choiceTF
{Mọi nhiệt lượng của dây nung đều chắc chắn chỉ dùng để nóng chảy mẫu.}
{\True Khi giải dạng “Thực nghiệm xác định nhiệt nóng chảy riêng”, cần xác định đúng trạng thái đầu, trạng thái cuối và quy ước dấu hoặc điều kiện áp dụng.}
{Có thể bỏ qua mọi dữ kiện về khối lượng, nhiệt độ và hiệu suất trong mọi bài toán thuộc dạng này.}
{\True Có thể xác định λ từ năng lượng cung cấp và khối lượng nóng chảy do nguồn.}
\loigiai{
\begin{itemchoice}
\item (Sai) Mọi nhiệt lượng của dây nung đều chắc chắn chỉ dùng để nóng chảy mẫu. (Vì): Mọi nhiệt lượng của dây nung đều chắc chắn chỉ dùng để nóng chảy mẫu. b) Đ: Khi giải dạng “Thực nghiệm xác định nhiệt nóng chảy riêng”, cần xác định đúng trạng thái đầu, trạng thái cuối và quy ước dấu hoặc điều kiện áp dụng. c) S: Có thể bỏ qua mọi dữ kiện về khối lượng, nhiệt độ và hiệu suất trong mọi bài toán thuộc dạng này. d) Đ: Có thể xác định λ từ năng lượng cung cấp và khối lượng nóng chảy do nguồn. Kết luận dựa trên định nghĩa, điều kiện áp dụng và quy ước trong SGK KNTT
\item (Đúng) Khi giải dạng “Thực nghiệm xác định nhiệt nóng chảy riêng”, cần xác định đúng trạng thái đầu, trạng thái cuối và quy ước dấu hoặc điều kiện áp dụng.
\item (Sai) Có thể bỏ qua mọi dữ kiện về khối lượng, nhiệt độ và hiệu suất trong mọi bài toán thuộc dạng này.
\item (Đúng) Có thể xác định λ từ năng lượng cung cấp và khối lượng nóng chảy do nguồn.
\end{itemchoice}
}
\end{ex}

% ===== Câu 107 =====
% ID: L12C1B6-204-TL
% Mức: VDC
\begin{bt}
% ID: L12C1B6-204-TL
% Mức: VDC
Trình bày cách nhận biết và phương pháp giải dạng “Thực nghiệm xác định nhiệt nóng chảy riêng”. Phân tích nhận định sau và sửa lại nếu sai: “Có thể xác định λ từ năng lượng cung cấp và khối lượng nóng chảy do nguồn.”
\loigiai{
Trước hết xác định hiện tượng, trạng thái và các đại lượng đã cho. Nêu định nghĩa hoặc hệ thức phù hợp, đổi về đơn vị SI, áp dụng đúng điều kiện và kiểm tra ý nghĩa kết quả. Nhận định cần đối chiếu: Có thể xác định λ từ năng lượng cung cấp và khối lượng nóng chảy do nguồn. Vì vậy nhận định đã cho đúng.
}
\end{bt}

% ===== Câu 110 =====
% ID: L12C1B6-205-DS
% Mức: VD
\begin{ex}
% ID: L12C1B6-205-DS
% Mức: VD
Xét các phát biểu sau về dạng “Thực nghiệm xác định nhiệt nóng chảy riêng” (bộ số 2).
\choiceTF
{\True Khi giải dạng “Thực nghiệm xác định nhiệt nóng chảy riêng”, cần xác định đúng trạng thái đầu, trạng thái cuối và quy ước dấu hoặc điều kiện áp dụng.}
{Có thể bỏ qua mọi dữ kiện về khối lượng, nhiệt độ và hiệu suất trong mọi bài toán thuộc dạng này.}
{\True Có thể xác định λ từ năng lượng cung cấp và khối lượng nóng chảy do nguồn.}
{Mọi nhiệt lượng của dây nung đều chắc chắn chỉ dùng để nóng chảy mẫu.}
\loigiai{
\begin{itemchoice}
\item (Đúng) Khi giải dạng “Thực nghiệm xác định nhiệt nóng chảy riêng”, cần xác định đúng trạng thái đầu, trạng thái cuối và quy ước dấu hoặc điều kiện áp dụng. (Vì): Khi giải dạng “Thực nghiệm xác định nhiệt nóng chảy riêng”, cần xác định đúng trạng thái đầu, trạng thái cuối và quy ước dấu hoặc điều kiện áp dụng. b) S: Có thể bỏ qua mọi dữ kiện về khối lượng, nhiệt độ và hiệu suất trong mọi bài toán thuộc dạng này. c) Đ: Có thể xác định λ từ năng lượng cung cấp và khối lượng nóng chảy do nguồn. d) S: Mọi nhiệt lượng của dây nung đều chắc chắn chỉ dùng để nóng chảy mẫu. Kết luận dựa trên định nghĩa, điều kiện áp dụng và quy ước trong SGK KNTT
\item (Sai) Có thể bỏ qua mọi dữ kiện về khối lượng, nhiệt độ và hiệu suất trong mọi bài toán thuộc dạng này.
\item (Đúng) Có thể xác định λ từ năng lượng cung cấp và khối lượng nóng chảy do nguồn.
\item (Sai) Mọi nhiệt lượng của dây nung đều chắc chắn chỉ dùng để nóng chảy mẫu.
\end{itemchoice}
}
\end{ex}

% ===== Câu 114 =====
% ID: L12C1B6-206-DS
% Mức: VDC
\begin{ex}
% ID: L12C1B6-206-DS
% Mức: VDC
Xét các phát biểu sau về dạng “Thực nghiệm xác định nhiệt nóng chảy riêng” (bộ số 3).
\choiceTF
{Có thể bỏ qua mọi dữ kiện về khối lượng, nhiệt độ và hiệu suất trong mọi bài toán thuộc dạng này.}
{\True Có thể xác định λ từ năng lượng cung cấp và khối lượng nóng chảy do nguồn.}
{Mọi nhiệt lượng của dây nung đều chắc chắn chỉ dùng để nóng chảy mẫu.}
{\True Khi giải dạng “Thực nghiệm xác định nhiệt nóng chảy riêng”, cần xác định đúng trạng thái đầu, trạng thái cuối và quy ước dấu hoặc điều kiện áp dụng.}
\loigiai{
\begin{itemchoice}
\item (Sai) Có thể bỏ qua mọi dữ kiện về khối lượng, nhiệt độ và hiệu suất trong mọi bài toán thuộc dạng này. (Vì): Có thể bỏ qua mọi dữ kiện về khối lượng, nhiệt độ và hiệu suất trong mọi bài toán thuộc dạng này. b) Đ: Có thể xác định λ từ năng lượng cung cấp và khối lượng nóng chảy do nguồn. c) S: Mọi nhiệt lượng của dây nung đều chắc chắn chỉ dùng để nóng chảy mẫu. d) Đ: Khi giải dạng “Thực nghiệm xác định nhiệt nóng chảy riêng”, cần xác định đúng trạng thái đầu, trạng thái cuối và quy ước dấu hoặc điều kiện áp dụng. Kết luận dựa trên định nghĩa, điều kiện áp dụng và quy ước trong SGK KNTT
\item (Đúng) Có thể xác định λ từ năng lượng cung cấp và khối lượng nóng chảy do nguồn.
\item (Sai) Mọi nhiệt lượng của dây nung đều chắc chắn chỉ dùng để nóng chảy mẫu.
\item (Đúng) Khi giải dạng “Thực nghiệm xác định nhiệt nóng chảy riêng”, cần xác định đúng trạng thái đầu, trạng thái cuối và quy ước dấu hoặc điều kiện áp dụng.
\end{itemchoice}
}
\end{ex}

% ===== Câu 233 =====
% ID: L12C1B6-207-TN
% Mức: NB
\begin{ex}
% ID: L12C1B6-207-TN
% Mức: NB
Câu 1 thuộc dạng “Thực nghiệm xác định nhiệt nóng chảy riêng”. Phát biểu nào sau đây đúng?
\choice
{Mọi nhiệt lượng của dây nung đều chắc chắn chỉ dùng để nóng chảy mẫu.}
{Nhiệt lượng và nhiệt độ là hai đại lượng hoàn toàn giống nhau.}
{Mọi quá trình nhiệt học đều không phụ thuộc điều kiện thực hiện.}
{\True Có thể xác định λ từ năng lượng cung cấp và khối lượng nóng chảy do nguồn.}
\loigiai{
Phát biểu đúng là: Có thể xác định λ từ năng lượng cung cấp và khối lượng nóng chảy do nguồn. Các phương án còn lại trái với mô hình, định nghĩa hoặc hệ thức nhiệt học tương ứng.
}
\end{ex}

% ===== Câu 235 =====
% ID: L12C1B6-208-TN
% Mức: NB
\begin{ex}
% ID: L12C1B6-208-TN
% Mức: NB
Câu 5 thuộc dạng “Thực nghiệm xác định nhiệt nóng chảy riêng”. Phát biểu nào sau đây đúng?
\choice
{Mọi nhiệt lượng của dây nung đều chắc chắn chỉ dùng để nóng chảy mẫu.}
{Nhiệt lượng và nhiệt độ là hai đại lượng hoàn toàn giống nhau.}
{Mọi quá trình nhiệt học đều không phụ thuộc điều kiện thực hiện.}
{\True Có thể xác định λ từ năng lượng cung cấp và khối lượng nóng chảy do nguồn.}
\loigiai{
Phát biểu đúng là: Có thể xác định λ từ năng lượng cung cấp và khối lượng nóng chảy do nguồn. Các phương án còn lại trái với mô hình, định nghĩa hoặc hệ thức nhiệt học tương ứng.
}
\end{ex}

% ===== Câu 237 =====
% ID: L12C1B6-209-TN
% Mức: NB
\begin{ex}
% ID: L12C1B6-209-TN
% Mức: NB
Câu 9 thuộc dạng “Thực nghiệm xác định nhiệt nóng chảy riêng”. Phát biểu nào sau đây đúng?
\choice
{Mọi nhiệt lượng của dây nung đều chắc chắn chỉ dùng để nóng chảy mẫu.}
{Nhiệt lượng và nhiệt độ là hai đại lượng hoàn toàn giống nhau.}
{Mọi quá trình nhiệt học đều không phụ thuộc điều kiện thực hiện.}
{\True Có thể xác định λ từ năng lượng cung cấp và khối lượng nóng chảy do nguồn.}
\loigiai{
Phát biểu đúng là: Có thể xác định λ từ năng lượng cung cấp và khối lượng nóng chảy do nguồn. Các phương án còn lại trái với mô hình, định nghĩa hoặc hệ thức nhiệt học tương ứng.
}
\end{ex}

% ===== Câu 239 =====
% ID: L12C1B6-210-TN
% Mức: TH
\begin{ex}
% ID: L12C1B6-210-TN
% Mức: TH
Câu 2 thuộc dạng “Thực nghiệm xác định nhiệt nóng chảy riêng”. Phát biểu nào sau đây đúng?
\choice
{Nhiệt lượng và nhiệt độ là hai đại lượng hoàn toàn giống nhau.}
{Mọi quá trình nhiệt học đều không phụ thuộc điều kiện thực hiện.}
{\True Có thể xác định λ từ năng lượng cung cấp và khối lượng nóng chảy do nguồn.}
{Mọi nhiệt lượng của dây nung đều chắc chắn chỉ dùng để nóng chảy mẫu.}
\loigiai{
Phát biểu đúng là: Có thể xác định λ từ năng lượng cung cấp và khối lượng nóng chảy do nguồn. Các phương án còn lại trái với mô hình, định nghĩa hoặc hệ thức nhiệt học tương ứng.
}
\end{ex}

% ===== Câu 241 =====
% ID: L12C1B6-211-TN
% Mức: TH
\begin{ex}
% ID: L12C1B6-211-TN
% Mức: TH
Câu 6 thuộc dạng “Thực nghiệm xác định nhiệt nóng chảy riêng”. Phát biểu nào sau đây đúng?
\choice
{Nhiệt lượng và nhiệt độ là hai đại lượng hoàn toàn giống nhau.}
{Mọi quá trình nhiệt học đều không phụ thuộc điều kiện thực hiện.}
{\True Có thể xác định λ từ năng lượng cung cấp và khối lượng nóng chảy do nguồn.}
{Mọi nhiệt lượng của dây nung đều chắc chắn chỉ dùng để nóng chảy mẫu.}
\loigiai{
Phát biểu đúng là: Có thể xác định λ từ năng lượng cung cấp và khối lượng nóng chảy do nguồn. Các phương án còn lại trái với mô hình, định nghĩa hoặc hệ thức nhiệt học tương ứng.
}
\end{ex}

% ===== Câu 243 =====
% ID: L12C1B6-212-TN
% Mức: TH
\begin{ex}
% ID: L12C1B6-212-TN
% Mức: TH
Câu 10 thuộc dạng “Thực nghiệm xác định nhiệt nóng chảy riêng”. Phát biểu nào sau đây đúng?
\choice
{Nhiệt lượng và nhiệt độ là hai đại lượng hoàn toàn giống nhau.}
{Mọi quá trình nhiệt học đều không phụ thuộc điều kiện thực hiện.}
{\True Có thể xác định λ từ năng lượng cung cấp và khối lượng nóng chảy do nguồn.}
{Mọi nhiệt lượng của dây nung đều chắc chắn chỉ dùng để nóng chảy mẫu.}
\loigiai{
Phát biểu đúng là: Có thể xác định λ từ năng lượng cung cấp và khối lượng nóng chảy do nguồn. Các phương án còn lại trái với mô hình, định nghĩa hoặc hệ thức nhiệt học tương ứng.
}
\end{ex}

% ===== Câu 245 =====
% ID: L12C1B6-213-TN
% Mức: VD
\begin{ex}
% ID: L12C1B6-213-TN
% Mức: VD
Câu 3 thuộc dạng “Thực nghiệm xác định nhiệt nóng chảy riêng”. Phát biểu nào sau đây đúng?
\choice
{Mọi quá trình nhiệt học đều không phụ thuộc điều kiện thực hiện.}
{\True Có thể xác định λ từ năng lượng cung cấp và khối lượng nóng chảy do nguồn.}
{Mọi nhiệt lượng của dây nung đều chắc chắn chỉ dùng để nóng chảy mẫu.}
{Nhiệt lượng và nhiệt độ là hai đại lượng hoàn toàn giống nhau.}
\loigiai{
Phát biểu đúng là: Có thể xác định λ từ năng lượng cung cấp và khối lượng nóng chảy do nguồn. Các phương án còn lại trái với mô hình, định nghĩa hoặc hệ thức nhiệt học tương ứng.
}
\end{ex}

% ===== Câu 247 =====
% ID: L12C1B6-214-TN
% Mức: VD
\begin{ex}
% ID: L12C1B6-214-TN
% Mức: VD
Câu 7 thuộc dạng “Thực nghiệm xác định nhiệt nóng chảy riêng”. Phát biểu nào sau đây đúng?
\choice
{Mọi quá trình nhiệt học đều không phụ thuộc điều kiện thực hiện.}
{\True Có thể xác định λ từ năng lượng cung cấp và khối lượng nóng chảy do nguồn.}
{Mọi nhiệt lượng của dây nung đều chắc chắn chỉ dùng để nóng chảy mẫu.}
{Nhiệt lượng và nhiệt độ là hai đại lượng hoàn toàn giống nhau.}
\loigiai{
Phát biểu đúng là: Có thể xác định λ từ năng lượng cung cấp và khối lượng nóng chảy do nguồn. Các phương án còn lại trái với mô hình, định nghĩa hoặc hệ thức nhiệt học tương ứng.
}
\end{ex}

% ===== Câu 249 =====
% ID: L12C1B6-215-TN
% Mức: VDC
\begin{ex}
% ID: L12C1B6-215-TN
% Mức: VDC
Câu 4 thuộc dạng “Thực nghiệm xác định nhiệt nóng chảy riêng”. Phát biểu nào sau đây đúng?
\choice
{\True Có thể xác định λ từ năng lượng cung cấp và khối lượng nóng chảy do nguồn.}
{Mọi nhiệt lượng của dây nung đều chắc chắn chỉ dùng để nóng chảy mẫu.}
{Nhiệt lượng và nhiệt độ là hai đại lượng hoàn toàn giống nhau.}
{Mọi quá trình nhiệt học đều không phụ thuộc điều kiện thực hiện.}
\loigiai{
Phát biểu đúng là: Có thể xác định λ từ năng lượng cung cấp và khối lượng nóng chảy do nguồn. Các phương án còn lại trái với mô hình, định nghĩa hoặc hệ thức nhiệt học tương ứng.
}
\end{ex}

% ===== Câu 251 =====
% ID: L12C1B6-216-TN
% Mức: VDC
\begin{ex}
% ID: L12C1B6-216-TN
% Mức: VDC
Câu 8 thuộc dạng “Thực nghiệm xác định nhiệt nóng chảy riêng”. Phát biểu nào sau đây đúng?
\choice
{\True Có thể xác định λ từ năng lượng cung cấp và khối lượng nóng chảy do nguồn.}
{Mọi nhiệt lượng của dây nung đều chắc chắn chỉ dùng để nóng chảy mẫu.}
{Nhiệt lượng và nhiệt độ là hai đại lượng hoàn toàn giống nhau.}
{Mọi quá trình nhiệt học đều không phụ thuộc điều kiện thực hiện.}
\loigiai{
Phát biểu đúng là: Có thể xác định λ từ năng lượng cung cấp và khối lượng nóng chảy do nguồn. Các phương án còn lại trái với mô hình, định nghĩa hoặc hệ thức nhiệt học tương ứng.
}
\end{ex}

% ===== Câu 266 =====
% ID: L12C1B6-217-TLN
% Mức: NB
\begin{ex}
% ID: L12C1B6-217-TLN
% Mức: NB
Cho bốn nhận định: (1) Có thể xác định λ từ năng lượng cung cấp và khối lượng nóng chảy do nguồn. (2) Mọi nhiệt lượng của dây nung đều chắc chắn chỉ dùng để nóng chảy mẫu. (3) Cần kiểm tra đơn vị và điều kiện áp dụng khi xử lí dạng “Thực nghiệm xác định nhiệt nóng chảy riêng”. (4) Có thể thay tùy ý nhiệt độ Celsius cho nhiệt độ tuyệt đối trong mọi công thức. Có bao nhiêu nhận định đúng? Trả lời bằng một số nguyên.
\shortans{2}
\loigiai{
Nhận định (1) và (3) đúng; (2) và (4) sai. Vậy có 2 nhận định đúng.
}
\end{ex}

% ===== Câu 268 =====
% ID: L12C1B6-218-TLN
% Mức: TH
\begin{ex}
% ID: L12C1B6-218-TLN
% Mức: TH
Cho bốn nhận định: (1) Có thể xác định λ từ năng lượng cung cấp và khối lượng nóng chảy do nguồn. (2) Mọi nhiệt lượng của dây nung đều chắc chắn chỉ dùng để nóng chảy mẫu. (3) Cần kiểm tra đơn vị và điều kiện áp dụng khi xử lí dạng “Thực nghiệm xác định nhiệt nóng chảy riêng”. (4) Có thể thay tùy ý nhiệt độ Celsius cho nhiệt độ tuyệt đối trong mọi công thức. Có bao nhiêu nhận định đúng? Trả lời bằng một số nguyên.
\shortans{2}
\loigiai{
Nhận định (1) và (3) đúng; (2) và (4) sai. Vậy có 2 nhận định đúng.
}
\end{ex}

% ===== Câu 270 =====
% ID: L12C1B6-219-TLN
% Mức: VD
\begin{ex}
% ID: L12C1B6-219-TLN
% Mức: VD
Cho bốn nhận định: (1) Có thể xác định λ từ năng lượng cung cấp và khối lượng nóng chảy do nguồn. (2) Mọi nhiệt lượng của dây nung đều chắc chắn chỉ dùng để nóng chảy mẫu. (3) Cần kiểm tra đơn vị và điều kiện áp dụng khi xử lí dạng “Thực nghiệm xác định nhiệt nóng chảy riêng”. (4) Có thể thay tùy ý nhiệt độ Celsius cho nhiệt độ tuyệt đối trong mọi công thức. Có bao nhiêu nhận định đúng? Trả lời bằng một số nguyên.
\shortans{2}
\loigiai{
Nhận định (1) và (3) đúng; (2) và (4) sai. Vậy có 2 nhận định đúng.
}
\end{ex}

% ===== Câu 272 =====
% ID: L12C1B6-220-TLN
% Mức: VD
\begin{ex}
% ID: L12C1B6-220-TLN
% Mức: VD
Cho bốn nhận định: (1) Có thể xác định λ từ năng lượng cung cấp và khối lượng nóng chảy do nguồn. (2) Mọi nhiệt lượng của dây nung đều chắc chắn chỉ dùng để nóng chảy mẫu. (3) Cần kiểm tra đơn vị và điều kiện áp dụng khi xử lí dạng “Thực nghiệm xác định nhiệt nóng chảy riêng”. (4) Có thể thay tùy ý nhiệt độ Celsius cho nhiệt độ tuyệt đối trong mọi công thức. Có bao nhiêu nhận định đúng? Trả lời bằng một số nguyên.
\shortans{2}
\loigiai{
Nhận định (1) và (3) đúng; (2) và (4) sai. Vậy có 2 nhận định đúng.
}
\end{ex}

% ===== Câu 274 =====
% ID: L12C1B6-221-TLN
% Mức: VDC
\begin{ex}
% ID: L12C1B6-221-TLN
% Mức: VDC
Cho bốn nhận định: (1) Có thể xác định λ từ năng lượng cung cấp và khối lượng nóng chảy do nguồn. (2) Mọi nhiệt lượng của dây nung đều chắc chắn chỉ dùng để nóng chảy mẫu. (3) Cần kiểm tra đơn vị và điều kiện áp dụng khi xử lí dạng “Thực nghiệm xác định nhiệt nóng chảy riêng”. (4) Có thể thay tùy ý nhiệt độ Celsius cho nhiệt độ tuyệt đối trong mọi công thức. Có bao nhiêu nhận định đúng? Trả lời bằng một số nguyên.
\shortans{2}
\loigiai{
Nhận định (1) và (3) đúng; (2) và (4) sai. Vậy có 2 nhận định đúng.
}
\end{ex}

% ===== Câu 276 =====
% ID: L12C1B6-222-TLN
% Mức: VDC
\begin{ex}
% ID: L12C1B6-222-TLN
% Mức: VDC
Cho bốn nhận định: (1) Có thể xác định λ từ năng lượng cung cấp và khối lượng nóng chảy do nguồn. (2) Mọi nhiệt lượng của dây nung đều chắc chắn chỉ dùng để nóng chảy mẫu. (3) Cần kiểm tra đơn vị và điều kiện áp dụng khi xử lí dạng “Thực nghiệm xác định nhiệt nóng chảy riêng”. (4) Có thể thay tùy ý nhiệt độ Celsius cho nhiệt độ tuyệt đối trong mọi công thức. Có bao nhiêu nhận định đúng? Trả lời bằng một số nguyên.
\shortans{2}
\loigiai{
Nhận định (1) và (3) đúng; (2) và (4) sai. Vậy có 2 nhận định đúng.
}
\end{ex}
